\chapter{Datos y variables}
\label{chap:datos}

% ==========================================================================
% BASE: capítulo 3 de la versión 2024 (pp. 12-16): panel código postal-mes,
% descriptivos, promedios por alcaldía, CP con más infracciones/incidentes.
% AJUSTES:
%  - Unidad: colonia (unidad territorial IECM), no código postal.
%  - Nueva fuente principal de outcome: víctimas FGJ (2019-jul 2024).
%  - Nueva construcción del tratamiento: exposición a radares por buffer.
%  - Documentar la auditoría de la muestra (Horacio: "mejorar construcción y
%    documentación de la muestra").
%  - Documentar los problemas de cobertura de infracciones (abr-jun 2021).
% ==========================================================================

\section{Fuentes}

El \emph{outcome} principal proviene de las carpetas de investigación de la
Fiscalía General de Justicia de la Ciudad de México (FGJ), que registran a
cada persona involucrada en un hecho de tránsito entre enero de 2019 y julio
de 2024. De este universo se conservan únicamente las lesiones y los
homicidios culposos por tránsito vehicular, lo que deja alrededor de
27{,}000 víctimas, de las cuales cerca de 4{,}400 fallecieron; el 88\% de los
registros cuenta con coordenadas propias, y el resto se ubica por coincidencia
de nombre de colonia. Como \emph{outcome} secundario se usa la serie
georreferenciada de incidentes viales reportados a los sistemas de atención de
emergencias del C5, disponible de forma continua desde 2014 hasta febrero de
2024; a diferencia de las víctimas, este registro no distingue gravedad ni
depende de que se abra una carpeta de investigación, por lo que sirve como
punto de comparación frente a un posible efecto de reporte
(Sección~\ref{sec:contexto_espacial}).

El mecanismo del diseño —si la reforma efectivamente incrementó la actividad
de fiscalización— se documenta con la base histórica de infracciones al
Reglamento de Tránsito, que cubre de enero de 2020 a abril de 2023 y en la que
el artículo 9 (exceso de velocidad) concentra por sí solo cerca del 77\% de
los registros. El tratamiento, por su parte, se construye a partir del
inventario anual de radares fijos y móviles y equipos asociados de
Fotocívicas, disponible de 2019 a 2024 y georreferenciado equipo por equipo,
según se describe en la sección siguiente. Todas estas fuentes se agregan
sobre la misma malla territorial: las 1{,}814 unidades del catálogo de
colonias del Instituto Electoral de la Ciudad de México, que es también la
unidad de observación del panel.

\section{Panel colonia-mes}

Se construye un panel balanceado de 1{,}814 colonias por 103 meses (enero de
2016 a julio de 2024). Cada \emph{outcome} se agrega al nivel colonia-mes por
asignación espacial (punto en polígono) o, para las infracciones sin
coordenadas a partir de mayo de 2021, por coincidencia de nombre de colonia.
Cada fuente lleva un indicador de cobertura temporal; fuera de su periodo de
cobertura los conteos se dejan como ausentes, no como cero.

\section{Construcción de la exposición a la infraestructura}
\label{sec:datos_exposicion}

Para cada colonia y cada año entre 2019 y 2024 se calcula, a partir del
inventario georreferenciado de radares fijos y móviles, tres medidas de
exposición: el número de equipos a menos de $b$ metros del polígono de la
colonia, la distancia al equipo más cercano, y la proporción del área de la
colonia que cae dentro de ese radio. Estas tres medidas se calculan para
cuatro definiciones de \emph{buffer} ($b \in \{250, 500, 1000, 1500\}$
metros), y de la primera de ellas se deriva el indicador binario
$\text{tratada}_i(b)$ que se usa como definición principal del tratamiento en
el capítulo~\ref{chap:estrategia}.

El diseño usa el inventario de 2021 —el más cercano y predeterminado respecto
al evento de abril de 2022— para fijar qué colonias cuentan como tratadas. La
tabla~\ref{tab:tratamiento_buffer} muestra que esta definición es, como es de
esperarse, sensible al radio elegido: pasa de 215 colonias tratadas (11.9\,\%
del total) con $b=250$\,m a 837 (46.1\,\%) con $b=1500$\,m. El capítulo de
resultados reporta el efecto para $b \in \{250, 500, 1000\}$ como análisis de
sensibilidad. Vale la pena notar que, aunque el número de radares por año
fluctúa moderadamente (Sección~\ref{sec:contexto_politica}), el conjunto de
colonias tratadas con $b=500$\,m se mantiene relativamente estable entre 2019
y 2024 (16.9\,\% a 20.2\,\%): la infraestructura se reubica poco de un año a
otro, lo que hace razonable fijar el inventario de 2021 como definición de
tratamiento predeterminada al evento.

\begin{table}[H]
  \centering
  \caption{Colonias tratadas por definición de \emph{buffer}, inventario 2021}
  \label{tab:tratamiento_buffer}
  \begin{tabular}{lrrr}
    \toprule
    Buffer & Colonias tratadas & Total & \% tratadas \\
    \midrule
    250 m  & 215 & 1{,}814 & 11.9\,\% \\
    500 m  & 328 & 1{,}814 & 18.1\,\% \\
    1{,}000 m & 590 & 1{,}814 & 32.5\,\% \\
    1{,}500 m & 837 & 1{,}814 & 46.1\,\% \\
    \bottomrule
  \end{tabular}
  \imagesource{elaboración propia con el inventario de radares y el catálogo
  de colonias del IECM.}
\end{table}

\section{Descriptivos}
\label{sec:datos_descriptivos}

La tabla~\ref{tab:descriptivos} resume la distribución de los tres
\emph{outcomes} y la variable de mecanismo a nivel colonia-mes, calculada
únicamente dentro del periodo de cobertura de cada fuente. El rasgo más
relevante para la estrategia empírica es la dispersión del \emph{outcome}
principal: las víctimas de tránsito promedian apenas 0.16 por colonia-mes, con
88.7\,\% de observaciones en cero, y los fallecidos son aún más raros
(97.9\,\% de ceros). Esta concentración en cero es la razón por la que el
capítulo~\ref{chap:resultados} contrasta el modelo lineal con un modelo de
conteo (Poisson): una variable así de dispersa se presta a que un puñado de
colonias con recuentos altos domine un ajuste por mínimos cuadrados, aun bajo
efectos fijos. Las infracciones, en contraste, son mucho menos dispersas en
proporción pero mucho más asimétricas en magnitud (máximo de 81{,}161 en una
sola colonia-mes), reflejo directo de la concentración espacial documentada
en la Sección~\ref{sec:contexto_espacial}.

\begin{table}[H]
  \centering
  \caption{Descriptivos a nivel colonia-mes}
  \label{tab:descriptivos}
  \begin{tabular}{lrrrrrr}
    \toprule
    Variable & Media & DE & P50 & P90 & Máx. & \% cero \\
    \midrule
    Incidentes C5        & 9.81   & 14.74   & 4 & 28 & 192    & 25.1\,\% \\
    Víctimas de tránsito  & 0.16   & 0.56    & 0 & 1  & 17     & 88.7\,\% \\
    \quad de las cuales, fallecidos & 0.03 & 0.21 & 0 & 0 & 17 & 97.9\,\% \\
    Infracciones          & 116.00 & 1{,}522.05 & 0 & 6 & 81{,}161 & 75.6\,\% \\
    \bottomrule
  \end{tabular}
  \imagesource{elaboración propia. Cada variable se calcula solo dentro del
  periodo de cobertura de su fuente (Sección~\ref{sec:datos_auditoria}).}
\end{table}

Antes de pasar al estudio de evento, vale la pena mirar simplemente los
niveles: en el periodo previo a abril de 2022, las colonias que resultan
tratadas con $b=500$\,m ya registraban, en promedio, casi el triple de
víctimas (0.39 vs. 0.14 por mes) y de incidentes C5 (18.4 vs. 6.9), y más de
doce veces las infracciones (528 vs. 43) que las colonias de control. Esta
brecha de niveles no es sorprendente —la infraestructura de fiscalización se
instala en vías primarias, que por definición concentran más tráfico y más
siniestralidad— pero es la razón por la que el capítulo~\ref{chap:estrategia}
identifica el efecto a partir de \emph{cambios} dentro de cada colonia
(efectos fijos) y no de una comparación de niveles entre tratadas y control.

\section{Auditoría de la muestra}
\label{sec:datos_auditoria}

Cada fuente se asigna a una colonia con un método distinto, y con una tasa de
éxito distinta. Las víctimas de tránsito se ubican por coordenadas en el
88\,\% de los casos; el 12\,\% restante, sin coordenadas válidas, se
descarta —una pérdida que no hay razón para suponer sistemáticamente
correlacionada con el tratamiento—. Los incidentes C5 llegan ya
georreferenciados de un procesamiento previo, con una tasa de asignación
cercana al 100\,\%. Las infracciones son la fuente más problemática: solo el
año 2020 trae coordenadas, con una asignación del 97\,\%; de 2021 en adelante
la asignación depende de hacer coincidir el nombre de colonia reportado con el
catálogo del IECM, y la tasa cae a alrededor de 55\,\% en 2022-2023 (65\,\%
en el conjunto del periodo). La causa es que el catálogo del IECM subdivide
varias colonias grandes en secciones —\texttt{GUERRERO I/II/III/IV},
\texttt{DOCTORES I-V}, nueve secciones distintas de \texttt{POLANCO}— que no
tienen contraparte en el nombre genérico que reportan las infracciones, y sin
coordenadas no hay forma de desambiguar a cuál sección pertenece cada
registro. Este problema afecta de forma desproporcionada a las colonias con
más infraestructura de fiscalización (Miguel Hidalgo, Álvaro Obregón), lo que
introduce un sesgo hacia el grupo tratado en cualquier estadístico agregado de
infracciones no asignadas. Es una razón adicional, junto con la ausencia de
outcome posterior para el evento de 2023, por la que las infracciones se usan
en esta tesis únicamente como evidencia de mecanismo y no como variable de
tratamiento ni como outcome principal.

La base de infracciones tiene además dos huecos temporales que se excluyen
explícitamente de su indicador de cobertura: abril de 2021, ausente por
completo de la fuente, y mayo-junio de 2021, en los que cerca del 67\,\% de
los registros del bimestre correspondiente (\texttt{2021\_b3}) llegan con la
colonia marcada como \texttt{DESCONOCIDO}. Ambos son problemas de la fuente
administrativa, no artefactos del procesamiento: se documentan aquí en vez de
imputarse. Por esta razón, y para mantener el análisis fuera de la fase más
aguda de la pandemia, la ventana del estudio de evento
(Sección~\ref{sec:estrategia_outcomes}) no arranca sino hasta junio de 2021,
diez meses antes del evento de abril de 2022.
